\section{Routing and Update Components}
\label{sec:routing}

\subsection{Three operational attention scores}

The state analysis asks what is represented. Attention is used only to
nominate components that may route information into those states. Let
\(A_{q,t}^{\ell h}\) be attention from registered query token \(q\) to token
\(t\) in layer \(\ell\), head \(h\). For the complete preregistered span
\(I_i\) of source occurrence \(i\), define
\begin{equation}
  a_i^{\ell h}(q)=\sum_{t\in I_i}A_{q,t}^{\ell h}.
\label{eq:source-mass}
\end{equation}
We sum over the complete span rather than selecting a favorable token after
inspection. Every report includes raw span mass, span-density
\(a_i^{\ell h}(q)/|I_i|\), and enrichment over a uniform-attention baseline,
so long source strings do not receive an unqualified advantage.

\paragraph{Broad aggregation candidate.}
At the Direct answer cue \(q_D\), let
\[
  M_G^{\ell h}=\sum_{i\in G_q}a_i^{\ell h}(q_D),
  \qquad
  r_i^{\ell h}=
  \frac{a_i^{\ell h}(q_D)}{M_G^{\ell h}+\epsilon}.
\]
For \(N\geq2\), define
\begin{equation}
  B_{\ell h}=
  \E_{\omega}\left[
  M_G^{\ell h}
  \frac{
    \exp\!\left(-\sum_{i\in G_q}r_i^{\ell h}
    \log(r_i^{\ell h}+\epsilon)\right)
  }{N}
  \right].
\label{eq:broad-score}
\end{equation}
The second factor is effective source coverage divided by \(N\): it is one
for uniform mass across all gold occurrences and \(1/N\) when mass is
concentrated on one. We also register the gold-versus-matched-negative
contrast
\begin{equation}
  C_{\ell h}^{B}=
  \E_\omega\left[
  \frac{1}{N}\sum_{i\in G_q}a_i^{\ell h}(q_D)
  -
  \frac{1}{|G_q^-|}\sum_{i\in G_q^-}a_i^{\ell h}(q_D)
  \right].
\label{eq:broad-contrast}
\end{equation}
A large \(B_{\ell h}\) without positive \(C_{\ell h}^B\) is nonspecific broad
access. A head or bundle earns the \emph{broad aggregation component} name
only if its query-local output intervention changes \(Z_D\) or the complete
count and a clean patch restores that endpoint.

\paragraph{Targeted retrieval candidate.}
At CoT retrieval step \(k\), the registered gold source \(g_k\) is defined by
\cref{eq:gold-next}; \(\mathcal V_{k-1}\) is the set already visited. Define
\begin{align}
  T_{\ell h}
  &=
  \E_\omega\E_{k\mid\omega}
  \left[
  \frac{a_{g_k}^{\ell h}(q_k^{\rm ret})}
       {\sum_{i=1}^{K}a_i^{\ell h}(q_k^{\rm ret})+\epsilon}
  \right],
\label{eq:target-score}\\
  D_{\ell h}
  &=
  \E_\omega\E_{k\mid\omega}
  \left[
  \frac{\sum_{i\in\mathcal V_{k-1}}a_i^{\ell h}(q_k^{\rm ret})}
       {\sum_{i=1}^{K}a_i^{\ell h}(q_k^{\rm ret})+\epsilon}
  \right].
\label{eq:duplicate-score}
\end{align}
We report raw \(a_{g_k}\), \(T_{\ell h}\), \(D_{\ell h}\), source-span top-1,
and the gold-versus-best-wrong QK margin. Attention to the source the model
actually emits, \(\pi_k\), is an execution-alignment statistic and never
replaces \(g_k\) in the candidate score. A targeted retrieval component must
have high gold-target mass, low visited-source mass, and a causal effect on
the emitted source identity or source-identity margin.

\paragraph{Progress-transition candidate.}
For a nonterminal controlled prefix, the next registered source is
\(g_{k+1}\). At the progress landmark \(q_k^{\rm end}\), define
\begin{equation}
  U_{\ell h}=
  \E_\omega\E_{k\mid\omega}
  \left[
  \frac{a_{g_{k+1}}^{\ell h}(q_k^{\rm end})}
       {\sum_{i=1}^{K}a_i^{\ell h}(q_k^{\rm end})+\epsilon}
  \right].
\label{eq:progress-score}
\end{equation}
Terminal prefixes are scored by the probability of \texttt{Total:}/stop and
by free-running stopping behavior rather than by forcing a nonexistent source
target. We use \emph{progress-transition component} or
\emph{next-source control}; ``successor head'' is reserved for a component
whose token-independent OV map survives relabeling and implements an abstract
\(k\mapsto k+1\) operation~\citep{gould2023successor}.

\begin{table}[H]
\centering
\small
\caption{Operational component roles. Scores nominate candidates; the final
column determines whether a functional name is licensed.}
\label{tab:head-roles}
\begin{tabularx}{\textwidth}{@{}P{.20\textwidth}P{.18\textwidth}P{.22\textwidth}Y@{}}
\toprule
Role & Registered query & Discovery evidence & Proximal causal endpoint \\
\midrule
Broad aggregation
  & Direct \(q_D\)
  & \(B_{\ell h}\), \(C_{\ell h}^B\), span-density sensitivity
  & \(Z_D\), complete numeric-string likelihood, exact count \\
Targeted retrieval
  & CoT \(q_k^{\rm ret}\)
  & \(T_{\ell h}\), low \(D_{\ell h}\), gold/wrong QK margin
  & gold/emitted source identity and unique coverage \\
Progress transition
  & CoT \(q_k^{\rm end}\)
  & \(U_{\ell h}\), terminal stop diagnostics
  & next source, repeat, continue, or stop \\
\bottomrule
\end{tabularx}
\end{table}

\subsection{Discovery, freezing, and family-equal aggregation}

Scores are computed at each eligible query, averaged over steps within a
rendered variant, then over variants within a family, and only then equally
over families. Discovery ranks candidates; validation freezes layer, bundle
size, score thresholds, and all causal endpoints. A family-wise max-statistic
permutation screen limits false discovery. We preregister at most five
candidates per role and model for the locked causal test and apply Holm
correction across the frozen role-by-endpoint tests, not across heatmaps
selected after inspection.

Attention maps are retained as interpretable summaries, but they do not
establish information flow. Ablations are applied to a head's contribution
after value aggregation and before/through output projection at the registered
query, not globally when a local intervention is available. Controls include
same-layer random heads, attention-mass-matched heads, wrong query rows,
lexically adjacent tokens, unrelated source spans, resample and mean ablation,
and clean-to-clean and corrupt-to-corrupt sham patches.

\begin{figure}[H]
\centering
\fbox{\parbox[c][6.0cm][c]{0.94\textwidth}{
\centering
\textbf{TODO Figure: Functional routing components}\\[0.5em]
\begin{minipage}[t]{.30\linewidth}\centering
\textit{Broad aggregation}\\
Direct answer-to-source map\\
\(B_{\ell h}\), contrast\\
\(Z_D\)/answer intervention
\end{minipage}\hfill
\begin{minipage}[t]{.30\linewidth}\centering
\textit{Targeted retrieval}\\
step-to-gold-source map\\
\(T_{\ell h}\), \(D_{\ell h}\)\\
identity/coverage intervention
\end{minipage}\hfill
\begin{minipage}[t]{.30\linewidth}\centering
\textit{Progress transition}\\
progress-to-next-source map\\
\(U_{\ell h}\), stop score\\
next/repeat/stop intervention
\end{minipage}\\[0.8em]
\small For each model and role, show a small number of real locked-test
heatmaps beside raw ablation damage, single-donor rescue, and matched controls.
}}
\caption{Planned component figure. Heatmaps describe candidate access
patterns; held-out head-output interventions determine function. Qwen and
Gemma head indices need not match.}
\label{fig:head-todo}
\end{figure}

\subsection{From heatmap to functional component}

Every role follows one certification ladder:
\begin{equation}
  \text{discover}
  \longrightarrow
  \text{freeze}
  \longrightarrow
  \text{held-out selectivity}
  \longrightarrow
  \text{local ablation}
  \longrightarrow
  \text{clean patch rescue}.
\label{eq:head-ladder}
\end{equation}
For an endpoint \(m\), normalized recovery
\[
  \operatorname{Rec}=
  \frac{m_{\rm patch}-m_{\rm corrupt}}
       {m_{\rm clean}-m_{\rm corrupt}}
\]
is reported only when the clean-corrupt denominator passes a validation-frozen
gate; all raw margins are always shown.

The intervention must first move the role's proximal semantic endpoint.
Broad candidates must affect terminal total quality. Targeted candidates must
first affect gold/emitted source identity or unique coverage.
Progress-transition candidates must first affect the next source, repetition,
or stopping. Pattern patches test localization; source-value and head-output
patches test information transport; residual patches localize no more finely
than the block. MLP output is tested separately after a block-level state
effect is established.

\begin{todobox}
\textbf{H1---frozen role screen.}
On discovery families, compute \(B_{\ell h}\), \(C_{\ell h}^B\),
\(T_{\ell h}\), \(D_{\ell h}\), and \(U_{\ell h}\) for every layer/head in
both models using step-to-variant-to-family aggregation. Freeze no more than
five candidates per role/model and all thresholds on validation families.
Use a family-wise max-statistic screen. On locked families, report raw span
mass, density and uniform-baseline sensitivity, gold-target versus
emitted-target alignment, score uncertainty, and representative heatmaps
without post-test reselection.
\end{todobox}

\begin{todobox}
\textbf{H2---functional component tests.}
For every frozen candidate, run query-local head-output ablation,
clean-to-corrupt single-donor patching, and path patching into the registered
state. Broad candidates must first affect \(Z_D\) or terminal total quality;
targeted candidates must first affect source identity/coverage; transition
candidates must first affect next source, repetition, or stop. Then measure
complete continuation likelihood and free-running exact count/MAE. Report
necessity, partial sufficiency, rescue, side effects, and every matched random,
same-layer, attention-mass, wrong-row, wrong-source, and sham control on the
same complete sequences.
\end{todobox}

\subsection{Cross-model functional replication}

Qwen3-8B and Gemma4-E4B share registered variables and causal gates, not head
numbers. The full Qwen analysis freezes the workflow; Gemma candidates are
then discovered on Gemma discovery families and evaluated on disjoint locked
families under the same role definitions. Replication means that the same
semantic intervention changes the same proximal and downstream variables even
if layers, head indices, or attention pictures differ. If only Qwen supports a
complete role, the claim is restricted and Gemma is reported as a boundary
case.
